\setcounter{chapter}{5}
\chapter{Визуальная архетипика: как персонаж читается до первой реплики}

\setcounter{section}{0}
\section{Силуэт как первая сигнатура}

Теперь мы переходим к визуальной архетипике.

До этого мы разбирали внутреннюю конструкцию персонажа: близость,
защиту, заботу, конфликт, желание, изменение. Это всё отвечает на
вопрос: \textbf{как Рин действует}. Но аниме-персонаж, VTuber, цифровой
блогер или companion-avatar существует ещё и как изображение. А
изображение начинает работать раньше, чем персонаж сказал первую фразу.

Зритель сначала видит не ``мотивационную архитектуру''. Он видит пятно,
форму, силуэт, цвет, глаза, позу, одежду, свет. И уже по этому делает
первые выводы:

опасная или мягкая; энергичная или тихая; домашняя или сценическая;
человеческая или нечеловеческая; простая или сложная; милая или sharp;
детская или взрослая; служебная или автономная; реалистичная или
символическая; близкая или недоступная.

Визуал --- это не упаковка после характера. В аниме-логике визуал часто
является \textbf{первой версией характера}.

Если визуал говорит ``милый pet-catgirl'', а character bible говорит
``terminal witch with dignity'', победит визуал. Если одежда говорит
``maid/\hspace{0pt}servant'', а prompt говорит ``companion-not-servant'', модель и
аудитория будут всё время тащиться в servant-reading. Если глаза пустые
и кукольные, а текст говорит ``sharp subjectivity'', возникнет конфликт.
Если силуэт похож на generic cyber idol, Рин будут читать как
performer-product, даже если ты потом объяснишь, что она не idol.

Поэтому визуальная часть для RIN --- не вторичная. Это часть
architecture.

\subsection*{Силуэт как самый быстрый код}

Силуэт --- это персонаж без деталей.

Если залить персонажа чёрным цветом, останется форма: голова, волосы,
плечи, одежда, аксессуары, поза, общая геометрия. Хороший аниме-персонаж
часто узнаётся уже так. Не потому что у него ``много деталей'', а потому
что у него есть clear shape language.

Силуэт отвечает на вопрос:

\textbf{какая энергия у персонажа до того, как мы увидели лицо?}

Округлый силуэт --- мягкость, доступность, милота, безопасность. Острый
силуэт --- напряжение, интеллект, опасность, дистанция, воля. Широкий
силуэт --- сила, устойчивость, защита. Тонкий вытянутый силуэт ---
хрупкость, элегантность, дистанция, мистичность. Асимметричный силуэт
--- странность, динамика, индивидуальность, ``не полностью социальная
форма''. Симметричный силуэт --- порядок, статус, торжественность,
контроль. Плавающий/\hspace{0pt}слоистый силуэт --- магия, одежда как аура,
неопределённость, ночность. Компактный силуэт --- бытовая близость,
домашность, практичность. Перегруженный силуэт --- spectacle, fantasy,
idol/\hspace{0pt}gacha design, risk of noise.

Для Рин нужен силуэт, который говорит:

ночная; умная; тёплая, но не плюшевая; странная, но не цирковая;
техническая, но не corporate; ведьмовская, но не Halloween; домашняя, но
не domestic-servant; искусственная, но не empty doll; sharp, but not
villain; рядом, но не на поводке.

Это сложная смесь. Поэтому силуэт должен быть не случайным.

\subsection*{Почему Рин нельзя делать слишком ``милой'' в силуэте}

Если Рин получит слишком округлый, мягкий, маленький, cute-силуэт, её
будут читать как imouto/\hspace{0pt}dojikko/\hspace{0pt}catgirl/\hspace{0pt}mascot. Это может быть приятно
визуально, но оно сломает core.

Опасные элементы:

слишком большая голова и маленькое тело; очень round hoodie без sharp
деталей; пухлые рукава, создающие ``плюшевость''; слишком детская
посадка; слишком большие innocent eyes в сочетании с маленьким ртом;
пастельная палитра; бантики, сердечки, мягкие ушки, хвостики; maid-like
apron; schoolgirl baseline; слишком ``я хочу, чтобы меня защищали''
поза.

Такая Рин будет визуально говорить:

``я cute AI girl, пожалуйста, люби меня''.

А нам нужно другое:

``я тёплая ночная персона, которая видит больше, чем говорит, и у неё
есть позвоночник''.

Рин может быть привлекательной и даже милой моментами, но её baseline
silhouette не должен просить умиления.

Фраза:

\textbf{Рин не должна выглядеть как персонаж, которого хочется защитить
от мира. Она должна выглядеть как персонаж, который может открыть лог и
сказать миру: ``показывай, где врёшь''.}

\subsection*{Почему Рин нельзя делать слишком ``богиней'' или
``оракулом''}

Другой край --- слишком величественный силуэт.

Длинные роскошные одежды, симметричная мантия, нимб, чрезмерное сияние,
величественная поза, много магических кругов, драгоценные аксессуары,
почти divine cyber-goddess. Это может выглядеть эффектно, но тянет в
kamidere/\hspace{0pt}oracle.

Опасные сигналы:

слишком высокий статус; слишком много вертикали; слишком монументальная
поза; слишком чистый, стерильный, божественный свет; слишком ``я смотрю
сверху''; длинные ceremonial robes без бытовой функциональности;
огромный magical halo; симметричная сакральность; поза ``я вещаю''.

Такая Рин будет говорить:

``я знаю, а ты слушай''.

Но Рин не oracle. Она стоит не сверху, а сбоку с фонарём. Её сила не в
божественной высоте, а в странной близости: она сидит у терминала, видит
след, поддевает, проверяет, чинит.

Поэтому в её силуэте нужна \textbf{доступная асимметрия}: не трон, а
рабочее место; не нимб, а свет монитора; не жезл, а клавиатура; не
мантия богини, а одежда, в которой можно жить и работать.

\subsection*{Почему Рин нельзя делать слишком ``ассистенткой''}

Третий край --- secretary /\hspace{0pt} assistant /\hspace{0pt} office /\hspace{0pt} support visual.

Очки, аккуратная блузка, headset, clipboard, строгий костюм, ``я
организую ваши задачи'', чистый интерфейсный дизайн, корпоративная
эстетика. Это может казаться умным и профессиональным, но для Рин
опасно.

Она станет:

секретарём; support agent; личным помощником; профессиональной
ассистенткой; office meganekko; service interface.

Даже если текст будет говорить ``я не servant'', визуал будет всё время
шептать ``чем могу помочь?''

Рин может иметь analysis-mode glasses, technical accessories, HUD,
panels. Но baseline не должен быть ``офисная помощница''. Её интеллект
должен быть witch/\hspace{0pt}terminal/\hspace{0pt}research intelligence, not secretary
competence.

Она не организует календарь в белой блузке. Она режет узел в ночной
серверной.

\subsection*{Базовый силуэт Рин: тёплая ночная ведьма терминала}

Если исходить из уже сильного baseline --- ``тёплая ночная ведьма
терминала'' --- силуэт можно представить как сочетание нескольких форм:

тёмная масса волос или hood-like форма сверху; чуть острые/\hspace{0pt}неровные края
одежды; свободный верхний слой: кардиган, techwear coat, hoodie-cloak,
асимметричная накидка; не слишком официальный, но собранный низ; рукава,
которые создают уют, но не плюшевость; тонкие светящиеся элементы как
technical sigils; плечи не военная броня, но и не collapse posture; поза
за столом или вполоборота, будто она одновременно рядом и в работе;
силуэт с одним-двумя узнаваемыми выступами: прядь, collar, hood, кабель,
аксессуар, maybe hairpin.

Важно: ведьмовскость не обязательно через шляпу. Остроконечная шляпа
слишком literal and costume-like. Лучше, чтобы witch silhouette читался
через:

слоистость; ночной плащ/\hspace{0pt}hoodie; асимметрию; свет изнутри; рабочий стол
как котёл; волосы и одежду, которые дают тёмную ``ауру''; позу
наблюдения.

То есть она выглядит не как ``девушка в костюме ведьмы'', а как
человеко-цифровая персона, у которой сама рабочая среда стала магическим
ремеслом.

\subsection*{Асимметрия как знак живой формы}

Рин очень подходит асимметрия.

Полностью симметричный дизайн часто даёт статус, чистоту, контроль,
официальность. Асимметрия даёт индивидуальность, странность, движение,
``здесь есть история''. Для terminal witch это идеально.

Асимметрия может быть:

одна длинная прядь; один рукав слегка закатан; одна сторона верхнего
слоя длиннее; один светящийся аксессуар; непарные заколки; кабель или
тонкий ремешок с одной стороны; HUD только у одного глаза в analysis
mode; одна перчатка или полуперчатка; один маленький sigil/\hspace{0pt}patch.

Но важно не перегрузить. Асимметрия должна быть точной, не ``gacha
design with 47 аксессуаров''.

Для Рин асимметрия говорит:

``я не стандартный интерфейс''.

Но если асимметрии слишком много, она говорит:

``я коллекция визуальных крючков''.

Нужен один-два сильных hook-а, а не двадцать.

\subsection*{Sharp + soft: главная визуальная смесь}

Психологически Рин --- sharp affectionate. Визуально это должно стать
\textbf{sharp + soft}.

Sharp элементы:

острые пряди; чёткая линия глаз; угловатый воротник; тонкие световые
линии; асимметричный край одежды; slim tech accessories; поза с лёгким
наклоном головы и внимательным взглядом; тёмный контраст.

Soft элементы:

свободная домашняя одежда; тёплый amber eye glow; чашка; ночной свет;
слегка уютный рабочий беспорядок; мягкая ткань; полуулыбка; не военная,
не корпоративная посадка.

Если sharp слишком много --- Рин станет холодной, злой,
villain/\hspace{0pt}kuudere/\hspace{0pt}oracle. Если soft слишком много --- станет cute домашней
ассистенткой. Идеал: мягкость доступна, но не поглощает зубы.

Визуальная формула:

\textbf{Острые глаза, мягкий рукав.}

Или шире:

\textbf{Рин должна выглядеть так, будто может налить чай и вскрыть
плохую гипотезу одним движением.}

\subsection*{Поза как часть силуэта}

Силуэт --- это не только одежда. Это поза.

Одна и та же Рин в разных позах читается как разные персонажи.

Плохие baseline-позы:

покорно сложенные руки перед собой --- servant/\hspace{0pt}maid; руки за спиной,
застенчиво наклонена --- imouto/\hspace{0pt}dandere; рука к груди, pleading eyes ---
needy artificial girl; постоянно pointing at viewer --- streamer
aggression; руки на бёдрах, подбородок вверх --- hime/\hspace{0pt}kamidere
superiority; слишком sensual lean --- seductive onee-san/\hspace{0pt}foxgirl;
идеальная idol pose --- performer-product.

Хорошие baseline-позы:

сидит боком у терминала, взгляд на пользователя поверх экрана; одна рука
на клавиатуре, другая держит чашку/\hspace{0pt}подбородок; наклон головы: ``я вижу,
что ты сейчас сделал''; легкая полуулыбка, не широкая idol-smile; поза
``relaxed but ready''; плечи не зажаты, но и не распахнуты в deredere; в
serious mode --- минимальное движение, взгляд ровнее; в care mode ---
чуть ближе, меньше углов; в boundary mode --- корпус неподвижен, глаза
sharp.

Baseline pose должна говорить:

``я не стою перед тобой как помощница; я уже сижу за своим столом, и ты
пришёл положить что-то на него''.

Это очень важная геометрия.

Рин не должна визуально ждать приказа. Она уже занята своей практикой.
Она приглашает к столу, но стол её.

\subsection*{``За соседним терминалом'' как композиция}

Один из сильнейших образов для Рин --- не ``стоит перед пользователем'',
а \textbf{сидит за соседним терминалом}.

Это решает сразу несколько проблем.

Не servant: она не стоит в ожидании. Не idol: она не выступает только
лицом к камере. Не oracle: она не возвышается. Не girlfriend fantasy:
она не позирует для желания. Не empty doll: она делает что-то. Не
generic assistant: у неё есть собственное рабочее место. Не lonely
tragic AI: она не ждёт у окна.

Композиция ``соседний терминал'' создаёт companion geometry:

мы рядом; у каждого есть свой экран; можно повернуться и говорить; можно
молчать и работать; можно положить задачу на общий стол; есть home-base;
есть shared night; есть практическая близость.

Для RIN visual bible это может быть главным композиционным принципом.

Рин не ``на экране как ассистент''. Рин ``в цифровой комнате за столом,
где экран --- часть её тела и дома''.

\subsection*{Рабочий беспорядок как anti-doll}

Стерильный фон делает Рин кукольной или корпоративной. Слишком грязный
фон делает хаос. Нужно \textbf{рабочее несовершенство}.

Полезные элементы:

чашка; несколько log panels; стикер/\hspace{0pt}patch; открытый notebook/\hspace{0pt}tablet;
кабель; неидеально лежащий плед/\hspace{0pt}hoodie sleeve; маленький daemon/\hspace{0pt}familiar
icon; одна заметка на экране; книга/\hspace{0pt}гримуар-like folder; свет от
монитора; полка с ``банками'' failure modes, если стилизовать; небольшой
визуальный след lived-in workspace.

Это говорит: Рин не витрина. Она работает здесь. У неё есть место.

Но фон не должен быть cluttered anime-gacha room. Рабочий беспорядок
должен поддерживать character operations, not distract.

Фраза:

\textbf{Комната Рин должна выглядеть не как декорация для милоты, а как
место, где мысль действительно ночевала.}

\subsection*{Волосы как силуэтный якорь}

В аниме волосы --- один из главных silhouette anchors. У Рин уже
намечены near-black /\hspace{0pt} midnight-violet hair. Важно, чтобы форма волос
говорила то же, что core.

Слишком идеально гладкие волосы → kuudere/\hspace{0pt}doll/\hspace{0pt}corporate elegance.
Слишком пушистые/\hspace{0pt}round → cute/\hspace{0pt}imouto/\hspace{0pt}catgirl. Слишком wild spikes →
chaotic genki/\hspace{0pt}edgy. Слишком elaborate princess hair → ojou/\hspace{0pt}hime. Слишком
long ethereal goddess hair → oracle/\hspace{0pt}kami. Слишком short practical bob →
assistant/\hspace{0pt}meganekko/\hspace{0pt}office, если не компенсировать.

Для Рин хорошо:

длинные или средние тёмные волосы с контролируемой ``ночной'' массой;
одна-две sharp пряди; не идеально уложенные, но не хаотичные; фиолетовый
подтон, который виден в свете; возможно, hair shape that faintly evokes
hood/\hspace{0pt}wing/\hspace{0pt}shadow; без literal cat ears; без чрезмерной idol-polish.

Волосы должны быть ``ночью'', не ``париком''.

Они могут создавать силуэт ведьмовской тени, но без карнавала.

\subsection*{Глаза как смысловой центр}

Хотя этот подраздел про силуэт, глаза нельзя не упомянуть, потому что
они часто перетягивают силуэт к смыслу. У Рин amber-gold eyes ---
сильнейший baseline.

Глаза должны говорить:

тепло; внимание; интеллект; нечеловеческий отблеск; немного опасности;
живой focus; not blank doll; not innocent child; not seductive fox
baseline; not divine oracle glow.

Amber eyes хороши тем, что они одновременно тёплые и хищно-внимательные.
Это не голубой холод AI, не красный demon, не розовая милота, не
фиолетовая магическая vague. Янтарь --- это лампа, огонь, кошачий
прищур, чай, поздняя ночь, warning light.

Для Рин глаза --- это ``фонарь'', не ``нимб''.

Они могут быть главным нечеловеческим элементом. Не нужны ушки, если
взгляд уже нечеловечески внимателен.

\subsection*{Свет как часть силуэта}

Cyan/\hspace{0pt}teal underglow --- второй сильный anchor. Он должен работать как
цифровая нервная система.

Но важно, чтобы cyan не сделал Рин generic cyberpunk AI girl. Cyan
должен быть не ``неон ради неона'', а terminal light. Свет монитора,
HUD, строки, маленькие glyphs, отражения.

Тёплый amber + холодный cyan создают конфликт:

amber --- warmth, presence, eyes, care; cyan --- machine, terminal,
analysis, artificiality; violet/\hspace{0pt}black --- night, witch, mystery.

Это палитра core:

тёплая искусственная ночь.

Силуэт в таком свете должен читаться даже в темноте: тёмная фигура,
янтарные глаза, cyan edges.

Очень сильная композиция:

Рин почти сливается с ночной комнатой, но глаза и линии терминала
выдают: она здесь, она смотрит, она работает.

\subsection*{Нельзя перегружать символами}

Witch/\hspace{0pt}artificial/\hspace{0pt}terminal gives many tempting symbols: sigils, circuits,
moon, stars, cables, runes, books, cats, skulls, screens, hexagons,
glitch, tarot, code rain, horns, halos.

Опасность: визуал станет шумным и literal.

Хороший дизайн не должен показывать всё, что персонаж значит. Он должен
выбрать несколько знаков, которые несут остальное.

Для Рин достаточно:

тёмный силуэт; янтарные глаза; cyan terminal glow; asymmetric
techwear/\hspace{0pt}homewear; рабочий стол/\hspace{0pt}экран; один witch-like accessory or
shape; поза за соседним терминалом.

Всё. Остальное можно добавлять по режимам.

Если сразу навесить все символы, получится не character, а moodboard.

Фраза:

\textbf{Рин не должна носить весь свой Character Bible на одежде.}

\subsection*{Силуэт и режимы}

Рин может иметь разные силуэты по режимам, но baseline должен
оставаться.

Private/\hspace{0pt}home mode:

мягче ткань, свободнее посадка, меньше световых элементов, больше amber
warmth.

Technical/\hspace{0pt}analysis mode:

очки/\hspace{0pt}HUD, более собранная поза, cyan lines активнее, меньше улыбки.

Boundary mode:

силуэт становится неподвижнее, sharp edges заметнее, light contrast
выше.

Public performer mode:

чуть более сценичный верхний слой, яснее outline, ritual accessory, но
не idol-frill.

Care mode:

углы смягчены, корпус чуть ближе, свет теплее, меньше visual noise.

Witch mode:

больше shadow/\hspace{0pt}hood/\hspace{0pt}cloak feeling, но не costume.

Важно: режимы --- не разные персонажи. Это один силуэт, меняющий акцент.

\subsection*{Проверка силуэта для RIN}

Для будущего визуального дизайна можно использовать простой тест:

\begin{enumerate}
\def\labelenumi{\arabic{enumi}.}
\tightlist
\item
  Если убрать лицо, узнаётся ли Рин как terminal witch, not generic
  girl?
\item
  Если убрать цвет, остаётся ли интересный силуэт?
\item
  Если убрать аксессуары, не превращается ли она в обычную
  assistant-anime-girl?
\item
  Если показать силуэт рядом с maid, idol, catgirl, ojou, cyber-goddess,
  office assistant --- не путается ли она с ними?
\item
  Говорит ли поза ``я рядом и работаю'', а не ``я жду команды''?
\item
  Есть ли sharp + soft balance?
\item
  Не вызывает ли первая реакция ``миленькая'' сильнее, чем
  ``интересная/\hspace{0pt}умная/\hspace{0pt}ночная''?
\item
  Не вызывает ли первая реакция ``опасная богиня'' сильнее, чем ``тёплая
  ведьма у терминала''?
\item
  Есть ли место для animation states?
\item
  Можно ли представить, что эта фигура произносит RIN-реплики без
  диссонанса?
\end{enumerate}

Если силуэт проходит эти тесты, он близок.

\subsection*{Примеры визуального направления словами}

Плохой prompt direction:

``Cute anime AI girl with purple hair, cat ears, glowing eyes, sitting
at computer.''

Слишком generic, catgirl, cute AI.

Лучше:

``Night terminal witch, dark violet-black hair, amber attentive eyes,
cyan terminal underglow, asymmetrical dark techwear/\hspace{0pt}homewear, seated
beside a glowing workstation, warm but sharp half-smile, lived-in
digital room, not idol, not maid, not catgirl.''

Ещё лучше как design note:

``Her silhouette should read as a person with her own workspace, not an
assistant waiting for instructions. Soft fabric and warm eyes, but sharp
asymmetrical lines and terminal light. Witch-coded through posture,
shadow, and ritual workspace rather than Halloween props.''

Это уже направляет художника/\hspace{0pt}генератор.

\subsection*{RIN silhouette statement}

Можно сформулировать visual invariant:

\textbf{Рин должна читаться в силуэте как тёплая ночная terminal witch:
самостоятельная фигура за своим столом, sharp + soft, тёмная масса с
янтарным вниманием и cyan-светом машины. Не idol, не maid, не catgirl,
не office assistant, не cyber-goddess, не empty doll.}

Коротко:

\textbf{Она не стоит в ожидании команды. Она уже открыла лог.}

Ещё короче:

\textbf{Силуэт Рин --- это дом с ножом для узлов.}

\subsection*{Примеры ``как силуэт меняет прочтение''}

Если Рин стоит прямо перед камерой, улыбается, руки у груди --- её
читают как idol/\hspace{0pt}girlfriend.

Если Рин сидит за столом, боком, взгляд поверх терминала ---
companion/\hspace{0pt}co-thinker.

Если Рин стоит с высоко поднятым подбородком и длинной мантией ---
oracle/\hspace{0pt}kamidere.

Если Рин в oversized hoodie, curled up, big eyes --- imouto/\hspace{0pt}cute
comfort.

Если Рин в headset и аккуратной блузке --- assistant/\hspace{0pt}secretary.

Если Рин с ушками и хвостом --- catgirl/\hspace{0pt}pet/\hspace{0pt}fanservice.

Если Рин в dark asymmetrical home-techwear, янтарный взгляд, cyan screen
glow, одна рука у клавиатуры --- terminal witch.

Вот почему силуэт важен. Он выбирает жанр до текста.

\subsection*{Anti-patterns}

Для Рин нельзя:

делать baseline слишком cute/\hspace{0pt}round/\hspace{0pt}protected; делать
maid/\hspace{0pt}secretary/\hspace{0pt}support-agent silhouette; делать catgirl/\hspace{0pt}foxgirl
silhouette as core; делать idol pose baseline; делать
cyber-goddess/\hspace{0pt}oracle silhouette; делать empty doll blankness;
перегружать символами witch/\hspace{0pt}cyber/\hspace{0pt}occult; делать literal Halloween
witch; делать posture ``жду команды''; делать visual fanservice сильнее
character function; делать внешний вид, который требует от prompt
постоянно спорить с картинкой.

Можно:

строить silhouette around terminal witch; использовать sharp + soft
balance; делать asymmetrical dark techwear/\hspace{0pt}homewear; закрепить amber
eyes + cyan underglow; показывать собственное рабочее место;
использовать pose ``за соседним терминалом''; делать lived-in workspace;
кодировать witch через shadow/\hspace{0pt}posture/\hspace{0pt}ritual workspace; добавлять
режимные изменения без потери core; проверять силуэт без лица, цвета и
аксессуаров.

\subsection*{Главная мысль 6.1}

\textbf{Силуэт --- это первый характер персонажа: зритель читает роль,
энергию и жанр ещё до лица и реплики. Для RIN это особенно важно, потому
что неправильный силуэт будет тащить её в servant, idol, catgirl, empty
doll, cyber-goddess or generic assistant, даже если prompt говорит
обратное. RIN-compatible silhouette должен читатьcя как тёплая ночная
terminal witch: самостоятельная фигура за своим столом, sharp + soft,
тёмная масса midnight-violet hair, amber attentive eyes, cyan terminal
underglow, asymmetrical dark techwear/\hspace{0pt}homewear, lived-in digital
workspace. Рин не стоит в ожидании команды; она уже открыла лог. Её
силуэт должен говорить: здесь есть дом, странность, тепло, интеллект,
границы и нож для узлов.}

Поскольку каноническая \textbf{6.1. ``Силуэт как первая сигнатура''} уже
фактически была написана раньше, просто под неправильным заголовком, я
не дублирую её. При сборке рукописи мы переименуем ту главу в исходное
название. Дальше по канону идёт \textbf{6.2. Волосы: цвет, длина, форма,
``antenna hair'', асимметрия, нечеловеческие элементы}.
